\section{Experiments and Results}
\label{sec:experiments}

\subsection{Setup}

We evaluate MLP distinguishers on three ciphers across the following round ranges: \speck (3--8), \simon (4--11), \present (2--7).
Each configuration uses 500K training samples, 3--5 random seeds, and the R2\_xor\_diff representation.
Compute budget: approximately 48 GPU-hours total on a single NVIDIA V100-PCIE-32GB.
Our experimental setting is 2-1-CT-R in the classification of~\citet{bellini2023dbitnet,gerault2024sok}.

\paragraph{Baselines.}
We compare against (i)~our best-bit-bias classical distinguisher, (ii)~Gohr's published ResNet numbers on \speck (92.4\% at 5r, 78.8\% at 6r, 61.6\% at 7r, 51.4\% at 8r)~\citep{gohr2019}, and (iii)~our own ResNet trained under identical conditions.

% ─────────────────────────────────────────────────────────────────────────
\subsection{Cross-Family Vulnerability}
\label{sec:exp-vulnerability}

\begin{table}[t]
\centering
\caption{Neural distinguisher accuracy (\%) vs.\ round count.
         Values are mean $\pm$ std over 3--5 seeds.
         95\% bootstrap CIs are reported for boundary rounds (accuracy $<$ 55\%).
         Dashes indicate untested configurations.}
\label{tab:baseline}
\begin{tabular}{@{}rccc@{}}
\toprule
\textbf{$r$} & \textbf{\speck (ARX)} & \textbf{\simon (Feistel)} & \textbf{\present (SPN)} \\
\midrule
2  & ---                    & ---                       & $100.00 \pm 0.00$ \\
3  & $99.99 \pm 0.00$       & ---                       & $100.00 \pm 0.00$ \\
4  & $97.78 \pm 0.42$       & $100.00 \pm 0.00$         & $98.39 \pm 0.03$ \\
5  & $87.37 \pm 0.49$       & $99.97 \pm 0.01$          & $71.92 \pm 0.23$ \\
6  & $66.96 \pm 0.73$       & $98.09 \pm 0.08$          & $52.55 \pm 0.39$ \\
7  & $50.76 \pm 0.10$       & $79.09 \pm 0.23$          & $50.03 \pm 0.02$ \\
8  & $49.98 \pm 0.20$       & $61.13 \pm 0.42$          & --- \\
9  & ---                    & $52.56 \pm 0.18$          & --- \\
10 & ---                    & $50.32 \pm 0.23$          & --- \\
11 & ---                    & $50.02 \pm 0.08$          & --- \\
\bottomrule
\end{tabular}
\end{table}

\begin{figure}[t]
  \centering
  \includegraphics[width=\columnwidth]{fig1_baseline_all_ciphers.pdf}
  \caption{Distinguisher accuracy vs.\ round count for all three cipher families.
  \simon (Feistel) remains above chance through 9 rounds; \speck (ARX) through 7; \present (SPN) through 6.
  The dashed line marks the 50\% random baseline.}
  \label{fig:baseline}
\end{figure}

\Cref{tab:baseline,fig:baseline} reveal a clear vulnerability hierarchy.
\simon remains distinguishable through 9 rounds (52.56\%, CI: [52.30, 52.71]), while \speck saturates at 7 rounds (50.76\%) and \present at 6 rounds (52.55\%, CI: [52.07, 53.04]).
As fractions of full cipher depth, the neural distinguisher penetrates 28\% of \simon, 32\% of \speck, and 19\% of \present.

The ordering is consistent with the structural properties of each family.
\simon's Feistel structure processes only half the state per round, requiring more rounds to achieve full diffusion.
\present's SPN design, with a bit permutation providing full diffusion in a single round, resists neural analysis most effectively.

% ─────────────────────────────────────────────────────────────────────────
\subsection{Neural vs.\ Classical Gap}
\label{sec:exp-classical}

\begin{table}[t]
\centering
\caption{Neural vs.\ classical (best-bit-bias) accuracy (\%) at selected rounds.
         Gap is in percentage points (pp).
         Gohr's published ResNet numbers are shown for reference on \speck.}
\label{tab:classical}
\begin{tabular}{@{}clcccc@{}}
\toprule
\textbf{Cipher} & \textbf{$r$} & \textbf{Classical} & \textbf{Neural (MLP)} & \textbf{Gap} & \textbf{Gohr ResNet} \\
\midrule
\multirow{4}{*}{\speck}
  & 3 & 75.00 & \textbf{99.99} & +25.0 & --- \\
  & 5 & 69.25 & \textbf{87.75} & +18.5 & 92.44 \\
  & 6 & 58.83 & \textbf{68.15} & +9.3  & 78.80 \\
  & 7 & 50.17 & 50.58          & +0.4  & 61.16 \\
\midrule
\multirow{4}{*}{\simon}
  & 4 & 75.11 & \textbf{100.00} & +24.9 & --- \\
  & 6 & 75.05 & \textbf{98.19}  & +23.1 & --- \\
  & 7 & 63.29 & \textbf{79.44}  & +16.2 & --- \\
  & 8 & 56.95 & \textbf{60.78}  & +3.8  & --- \\
\midrule
\multirow{4}{*}{\present}
  & 2 & 75.11 & \textbf{100.00} & +24.9 & --- \\
  & 4 & 61.34 & \textbf{98.27}  & +36.9 & --- \\
  & 5 & 54.24 & \textbf{72.32}  & +18.1 & --- \\
  & 6 & 51.23 & 52.38           & +1.2  & --- \\
\bottomrule
\end{tabular}
\end{table}

\begin{figure}[t]
  \centering
  \includegraphics[width=\columnwidth]{fig2_neural_vs_classical.pdf}
  \caption{Neural (MLP) vs.\ classical (best-bit-bias) accuracy across rounds.
  The neural advantage is largest at intermediate rounds and collapses at the boundary.}
  \label{fig:classical}
\end{figure}

\begin{figure}[t]
  \centering
  \includegraphics[width=\columnwidth]{fig4_neural_advantage.pdf}
  \caption{Neural advantage (neural $-$ classical, in pp) by cipher and round.
  \present at 4 rounds shows the largest gap (+36.9\pp).}
  \label{fig:advantage}
\end{figure}

The neural MLP consistently outperforms the best single-bit classical distinguisher by 4--37\pp in the exploitable range (accuracy $>$ 55\%), collapsing to $<$2\pp at boundary rounds (\Cref{tab:classical,fig:classical}).
\present at 4 rounds shows the largest gap (+36.9\pp, \Cref{fig:advantage}), suggesting that the SPN's high diffusion concentrates differential signal into multi-bit correlations invisible to single-bit analysis.

Note that our classical baseline is deliberately weak (single-bit).
Gohr's ResNet, shown for \speck, substantially outperforms our MLP: the ResNet achieves 61.2\% at 7 rounds where our MLP gets 50.8\%.
This gap is expected and intentional---our contribution is the cross-family comparison, not depth maximization.

% ─────────────────────────────────────────────────────────────────────────
\subsection{Transfer and Anti-Transfer}
\label{sec:exp-transfer}

\begin{table}[t]
\centering
\caption{Cross-round and cross-cipher transfer results.
         Source: \speck 5r, \simon 6r, \present 4r.
         $p$-values from one-sample $t$-test against 0.5.
         $\star$: $p<0.05$; $\star\star$: $p<0.005$.}
\label{tab:transfer}
\begin{tabular}{@{}llccl@{}}
\toprule
\textbf{Type} & \textbf{Configuration} & \textbf{Acc.\ (\%)} & \textbf{$p$-value} & \textbf{Verdict} \\
\midrule
\multicolumn{5}{l}{\textit{Cross-Cipher}} \\
  & SPECK$\to$SIMON   & $48.90 \pm 0.38$ & $0.004$      & $\star\star$ Anti \\
  & SIMON$\to$SPECK   & $49.92 \pm 0.07$ & $0.085$      & n.s. \\
\midrule
\multicolumn{5}{l}{\textit{Cross-Round: \speck, trained 5r}} \\
  & $\to$ 3r           & $45.53 \pm 0.06$ & $< 10^{-8}$  & $\star\star$ Anti \\
  & $\to$ 4r           & $45.49 \pm 0.09$ & $< 10^{-7}$  & $\star\star$ Anti \\
  & $\to$ 6r           & $50.47 \pm 0.36$ & $0.059$      & n.s. \\
\midrule
\multicolumn{5}{l}{\textit{Cross-Round: \simon, trained 6r}} \\
  & $\to$ 4r           & $48.61 \pm 0.10$ & $< 10^{-5}$  & $\star\star$ Anti \\
  & $\to$ 5r           & $48.83 \pm 0.34$ & $0.002$      & $\star\star$ Anti \\
  & $\to$ 7r           & $50.27 \pm 0.11$ & $0.009$      & $\star\star$ Weak pos. \\
\midrule
\multicolumn{5}{l}{\textit{Cross-Round: \present, trained 4r}} \\
  & $\to$ 2r           & $99.06 \pm 0.09$ & $< 10^{-11}$ & $\star\star$ Positive \\
  & $\to$ 3r           & $96.86 \pm 1.43$ & $< 10^{-6}$  & $\star\star$ Positive \\
  & $\to$ 5r           & $59.14 \pm 0.31$ & $< 10^{-6}$  & $\star\star$ Positive \\
  & $\to$ 6r           & $51.56 \pm 0.07$ & $< 10^{-6}$  & $\star\star$ Weak pos. \\
\bottomrule
\end{tabular}
\end{table}

\begin{figure}[t]
  \centering
  \includegraphics[width=\columnwidth]{fig3_transfer_comparison.pdf}
  \caption{Transfer polarity comparison.
  ARX (\speck) and Feistel (\simon) models score \emph{below} 50\% on lower-round data (anti-transfer),
  while SPN (\present) models score \emph{above} 50\% (positive transfer).
  Error bars: $\pm$1 std over 5 seeds.}
  \label{fig:transfer}
\end{figure}

This is the central result of the paper.
\Cref{tab:transfer,fig:transfer} show that transfer polarity depends on cipher family:

\paragraph{ARX (SPECK): anti-transfer.}
A 5-round \speck distinguisher evaluated on 3-round data scores 45.5\%---4.5\pp \emph{below} random chance ($p < 10^{-8}$).
The network has learned features that are not merely irrelevant at fewer rounds, but actively misleading.

\paragraph{Feistel (SIMON): anti-transfer.}
A 6-round \simon distinguisher evaluated on 4-round data scores 48.6\% ($p < 10^{-5}$).
The anti-transfer is weaker than SPECK (1.4\pp vs.\ 4.5\pp below chance), but statistically unambiguous.
This is unexpected: Gohr et al.~\citep{gohr2022assessment} proved that Feistel distinguishers learn only differential features, which should compose hierarchically.
We hypothesize that \simon's AND-rotation nonlinearity creates round-specific differential patterns that disrupt composition (\Cref{sec:discussion}).

\paragraph{SPN (PRESENT): positive transfer.}
A 4-round \present distinguisher scores 99.1\% on 2-round data and 96.9\% on 3-round data ($p < 10^{-6}$).
Features learned at higher rounds subsume those at lower rounds, consistent with differential-only feature learning.

\paragraph{Interpretation.}
The asymmetry between ARX/Feistel (anti-transfer) and SPN (positive transfer) provides direct empirical evidence about \emph{what} neural distinguishers learn.
SPN features are hierarchical: a model that can detect 4-round patterns trivially detects 2-round patterns.
ARX and Feistel features are round-specific: the patterns at 5 rounds not only differ from those at 3 rounds, they are \emph{negatively correlated} with the 3-round decision boundary.

% ─────────────────────────────────────────────────────────────────────────
\subsection{Key Recovery}
\label{sec:exp-key-recovery}

\begin{table}[t]
\centering
\caption{Bayesian key recovery results (5 seeds, 20K pairs per trial).
         Rank is out of 256 candidates per byte.}
\label{tab:key_recovery}
\begin{tabular}{@{}lcccccc@{}}
\toprule
\textbf{Cipher} & \textbf{Total $r$} & \textbf{Exact (\%)} & \textbf{Low rank} & \textbf{High rank} & \textbf{Top-10 (\%)} & \textbf{Dist.\ acc.} \\
\midrule
\speck (4r) & 4 & 40 & 1.0 & 1.6 & 100 & 100.0\% \\
\speck (5r) & 5 & 40 & 10.4 & 2.0 & 80  & 98.3\%  \\
\simon (5r) & 5 & 0  & 122.8 & 128.8 & 0 & 100.0\% \\
\bottomrule
\end{tabular}
\end{table}

Following Gohr's Bayesian scoring approach~\citep{gohr2019}, we train a distinguisher at $r{-}1$ rounds and search the 16-bit last-round subkey space in two phases (low byte, then high byte).
\speck key recovery succeeds: at 4 total rounds, the correct subkey ranks first in 100\% of trials; at 5 rounds, 40\% exact recovery with 80\% in the top 10 (\Cref{tab:key_recovery}).

\simon key recovery fails completely despite 100\% distinguisher accuracy.
Mean low-byte rank is 122.8/256---indistinguishable from random.
This negative result is structurally informative: \simon's Feistel round function applies the subkey through XOR with a nonlinear function of the state ($f(x) = (x \lll 1) \mathbin{\&} (x \lll 8) \oplus (x \lll 2)$), breaking the direct subkey--score relationship that Gohr's byte-by-byte search exploits in SPECK's modular addition.

% ─────────────────────────────────────────────────────────────────────────
\subsection{Ablation Summary}
\label{sec:exp-ablation}

\paragraph{Architecture independence (E13).}
Seven architectures trained on \speck at 5 rounds produce accuracies in [86.8\%, 88.7\%], a spread of 1.9\pp (\Cref{tab:arch-comparison}).
The MLP (88.7\%) slightly outperforms the CNN (87.9\%) and LSTM (87.6\%); the Siamese network (86.8\%) is lowest but still within statistical noise.
Full results in \Cref{sec:app-ablations}.

\begin{table}[t]
\centering
\caption{Architecture comparison on \speck 5r (mean accuracy over 5 seeds).}
\label{tab:arch-comparison}
\begin{tabular}{@{}lcccc@{}}
\toprule
\textbf{Architecture} & \textbf{Acc.\ (\%)} & \textbf{Params} & \textbf{Train (s)} & \textbf{Throughput} \\
\midrule
MLP           & $88.67 \pm 0.38$ & 191K  & 222  & 67.8K \\
CNN           & $87.92 \pm 0.43$ & 81K   & 288  & 86.2K \\
LSTM          & $87.57 \pm 0.44$ & 594K  & 273  & --- \\
GRU           & $87.64 \pm 0.37$ & 454K  & 250  & --- \\
Gohr-MLP      & $87.10 \pm 0.51$ & 457K  & 84   & 60.7K \\
Residual CNN  & $86.88 \pm 0.60$ & 135K  & 363  & 68.2K \\
Siamese       & $86.81 \pm 0.40$ & 57K   & 216  & 68.8K \\
\bottomrule
\end{tabular}
\end{table}

\paragraph{ResNet comparison (E13b).}
Gohr's depth-10 ResNet trained under our identical conditions improves over MLP by at most 2.64\pp (\present at 5r: 74.7\% vs.\ 72.0\%).
On \speck, the gap is $\leq$1.1\pp across all rounds; on \simon, $\leq$0.7\pp.
Notably, on \present at 3--4 rounds the MLP \emph{matches or exceeds} ResNet.
The vulnerability ordering (\simon $\gg$ \speck $>$ \present) is preserved under both architectures, confirming it is intrinsic to cipher structure (\Cref{tab:resnet}).

\begin{table}[t]
\centering
\caption{ResNet vs.\ MLP accuracy (\%) under identical training conditions.
         Gap = ResNet $-$ MLP.}
\label{tab:resnet}
\begin{tabular}{@{}clccc@{}}
\toprule
\textbf{Cipher} & \textbf{$r$} & \textbf{MLP} & \textbf{ResNet} & \textbf{Gap (pp)} \\
\midrule
\multirow{4}{*}{\speck}
  & 4 & 98.24 & 98.28 & +0.04 \\
  & 5 & 87.37 & 87.78 & +0.41 \\
  & 6 & 66.96 & 68.03 & +1.07 \\
  & 7 & 50.76 & 50.79 & +0.03 \\
\midrule
\multirow{4}{*}{\simon}
  & 5 & 99.95 & 99.95 & +0.00 \\
  & 6 & 98.09 & 98.21 & +0.12 \\
  & 7 & 79.00 & 79.29 & +0.29 \\
  & 8 & 60.86 & 61.53 & +0.67 \\
\midrule
\multirow{3}{*}{\present}
  & 3 & 100.00 & 99.99 & $-0.00$ \\
  & 4 & 98.45  & 98.25 & $-0.19$ \\
  & 5 & 72.01  & 74.65 & +2.64 \\
\bottomrule
\end{tabular}
\end{table}
